\documentclass{article}
\usepackage[utf8]{inputenc}
\usepackage[T1]{fontenc}
\usepackage{amsmath, amssymb}
\usepackage{fancyhdr}
\usepackage{geometry}
\usepackage[dvipsnames]{xcolor}
\usepackage{soul}
\usepackage{lmodern}
\usepackage{microtype}
\usepackage{amsthm}
\usepackage[hidelinks]{hyperref}


\newtheorem{theorem}{Theorem}[section]


\sethlcolor{yellow}

\geometry{a4paper, margin=2.5cm}

\title{The Congruence Lattice Problem: A Historical Survey}
\author{Štěpán Chrast}
\date{\today}

\pagestyle{fancy}
\fancyhf{}
\lhead{Štěpán Chrast}
\rhead{The Congruence Lattice Problem}
\cfoot{\thepage}
\setlength{\headheight}{14pt}

\begin{document}
\maketitle

\begin{abstract}
The Congruence Lattice Problem (CLP) asks whether every distributive algebraic lattice is isomorphic to the congruence lattice of a lattice. After early positive results in the finite case and several infinite cardinalities, the problem was solved negatively by Wehrung, who constructed a distributive algebraic lattice that is not representable as a congruence lattice. This article provides a historical and structural overview of the problem, the main positive results, and the ideas behind the counterexample and the optimal cardinal bound.
\end{abstract}
\tableofcontents
\newpage

\section{Introduction}

The Congruence Lattice Problem (CLP) asks whether every distributive algebraic lattice is
isomorphic to the congruence lattice $\operatorname{Con}(L)$ of some lattice $L$.
Here $\operatorname{Con}(L)$ denotes the lattice of all congruence relations of $L$,
ordered by inclusion (see Section~2 for definitions).
Traditionally attributed to Dilworth (1940s) and open for more than sixty years, the CLP became a central problem in the interaction between lattice theory and universal algebra.

The question is motivated by two general results about congruence lattices.
In 1942, Funayama and Nakayama proved that for every lattice $L$,
the lattice $\operatorname{Con}(L)$ is distributive \cite{FunayamaNakayama1942}.
In 1948, Birkhoff and Frink proved that for every algebra $A$,
$\operatorname{Con}(A)$ is an algebraic lattice \cite{BirkhoffFrink1948}.
(Definitions of distributive, compact, and algebraic are recalled in Section~2.)

Together, these observations yield the following fact about lattices:
\begin{quote}
\textit{For any lattice $L$, the congruence lattice of $L$ is a distributive algebraic lattice.}
\end{quote}

The CLP then asks whether these necessary conditions are also sufficient:
\begin{quote} \textit{Is every distributive algebraic lattice isomorphic to the congruence lattice of some lattice $L$?} \end{quote}

The aim of this paper is to present a brief historical and conceptual survey of
the CLP. After recalling the necessary background in lattice theory, we outline
the finite solution and the subsequent partial results for infinite lattices.
We then describe the structural obstructions that emerge in higher cardinalities
and show Wehrung’s counterexample, the negative resolution of the Congruence Lattice Problem,
and further results.

\section{Preliminaries}

We recall notions used throughout the paper:

\begin{description}
\item[\textbf{Distributive lattice.}] A lattice is  an algebra $(L,\vee,\wedge)$, where both $\vee$ and $\wedge$ satisfye idempotency, commutativity, associativity and absorption identities.
A lattice is called \textit{distributive} if it satisfies the distributive law
\[
x \wedge (y \vee z) = (x \wedge y) \vee (x \wedge z)
\]
for all $x,y,z \in L$.

\item[\textbf{$(\vee,0)$-semilattice.}]
A \emph{$(\vee,0)$-semilattice} is an algebra $S=(S,\vee,0)$ where $\vee$ is associative, commutative, and idempotent, and $0$ is a least element.

\item[\textbf{Congruence.}] A \textit{congruence relation} on a lattice $L$ is an equivalence relation $\theta$ such that: if $x \equiv_{\theta} y$ and $u \equiv_{\theta} v$, then
\[
x \vee u \equiv_{\theta} y \vee v, \qquad x \wedge u \equiv_{\theta} y \wedge v.
\]

\item[\textbf{Congruence lattice.}]
Let $L$ be a lattice. The set of all congruences of $L$, ordered by inclusion, forms a (complete) lattice denoted $\operatorname{Con}(L)$, called the \emph{congruence lattice} of $L$.

\item[\textbf{Principal congruence.}]
Let $L$ be a lattice and let $x,y\in L$. The least congruence $\theta$ of $L$ such that $x\equiv_\theta y$ is denoted by $\Theta_L(x,y)$ and is called the \emph{principal congruence} generated by $(x,y)$.

\item[\textbf{Complete lattice.}] A lattice is \textit{complete} if every subset has a supremum and an infimum.

\item[\textbf{Compact element.}] Let $L$ be a complete lattice. An element $x \in L$ is \textit{compact} if for every $A \subseteq L$,
\[
x \leq \bigvee A \;\Rightarrow\; \exists\, F_{finite} \subseteq A \text{ such that } x \leq \bigvee F .
\]

\item[\textbf{Algebraic lattice.}] A lattice $L$ is \textit{algebraic} if it is complete and every element is the supremum of the compact elements below it.

\item[\textbf{\boldmath$\operatorname{Con}_c(A)$.}]
Let $A$ be an algebra. The set of compact elements of the lattice $\operatorname{Con}(A)$ forms a $(\vee,0)$-semilattice, denoted by $\operatorname{Con}_c(A)$.

\item[\textbf{Representability.}] A distributive $(\vee,0)$-semilattice $S$ is \emph{representable} if
$S$ is isomorphic to the $\operatorname{Con}_c(L)$ for some lattice $L$, as $(\vee,0)$-semilattices.

For more details see \cite{GratzerGLT1998}.
\end{description}

\section{The Finite Case}

For the finite cases of lattices, the Congruence Lattice Problem was actually solved before the general CLP was proposed and the notion of algebraicity was introduced in this context.
In the 1940s, Dilworth discovered a theorem concerning finite distributive lattices, with a complete proof being later published by Grätzer and
Schmidt in 1963. The theorem states that:
\begin{theorem}
\textit{Every finite distributive lattice is isomorphic to the congruence lattice
of some finite lattice.}
\end{theorem}

For proof see \cite[Theorem~8.1, p.~101]{GratzerCFL}.
\\ \\
Since every finite lattice is complete and every element of a finite lattice is compact, every finite lattice is algebraic, this theorem provides a positive solution for finite lattices of the congruence lattice problem:
\begin{quote}
\textit{Every finite (and therefore algebraic) distributive lattice is isomorphic to the congruence lattice of some (finite) lattice.}
\end{quote}

\section{Partial Results for Infinite Lattices}

The finite solution of the Congruence Lattice Problem does not extend directly to infinite lattices, but some strong partial results were obtained for infinite distributive lattices.

In the infinite case, algebraicity is not automatic: not every distributive lattice is algebraic. Distributive lattices that are not algebraic are not isomorphic to the congruence lattice of some lattice, hence in the infinite setting we need to focus on algebraic distributive lattices.

A useful perspective is to work with compact congruences.
In the case of a lattice $L$, the compact elements of $\operatorname{Con}(L)$ form under join
a distributive $(\vee,0)$-semilattice, denoted $\operatorname{Con}_c(L)$.
This viewpoint yields an equivalent semilattice formulation of the Congruence Lattice Problem \cite{Wehrung2007AdvMath}:
\begin{quote}
\textit{Is every distributive $(\vee,0)$-semilattice $S$ isomorphic to $\operatorname{Con}_c(L)$ for some lattice $L$?}
\end{quote}

A sequence of positive results was obtained during the 1970s--1990s, establishing representability for broad classes of infinite algebraic distributive lattices and distributive $(\vee,0)$-semilattices.
Building on structural conditions introduced by E. T. Schmidt, several authors proved representability under certain cardinalities and other properties:

\begin{theorem}[Wehrung]
Every countable union of distributive lattices with zero is representable.
\end{theorem}

\begin{theorem}[Dobbertin]
Every distributive join-semilattice with zero in which every principal ideal is countable is representable.
\end{theorem}

\begin{theorem}[Huhn]
Every distributive join-semilattice with zero of cardinality at most $\aleph_1$ is representable.
\end{theorem}

for details, see \cite[p.~705]{Gratzer2007TwoProblems}.
\\ \\
Together, these and related results established representability for several broad classes of infinite distributive $(\vee,0)$-semilattices, covering cardinalities (on the semilattice side) up to $\aleph_1$. For a long time, no counterexample to the Congruence Lattice Problem was known, while partial positive solutions have been accumulating. Thus leading to further strengthening the confidence that the CLP could be true.

\section{Obstructions in the Infinite Case}

All those positive results show that many distributive $(\vee,0)$-semilattices
are representable at small cardinalities, but that does not automatically extend
to the general infinite case. In higher cardinalities, some structural obstructions emerge.

A key ingredient in Wehrung's proof is the notion of a \emph{weakly distributive}
homomorphism between join-semilattices \cite{Wehrung2007AdvMath}.
A $(\vee,0)$-homomorphism $\mu:S\to T$ is \emph{weakly distributive at $x\in S$} if,
whenever $\mu(x)\leq y_0\vee y_1$ in $T$, there exist $x_0,x_1\in S$ such that
$x\leq x_0\vee x_1$ and $\mu(x_i)\leq y_i$ for $i<2$.
We say that $\mu$ is \emph{weakly distributive} if it is weakly distributive at every $x\in S$
\cite{Wehrung2007AdvMath}.

In the construction of his counterexample, Wehrung considers homomorphisms
$\mu:\operatorname{Con}_c(A)\to S$ into a distributive $(\vee,0,1)$-semilattice $S$
and focuses on the situation where $1\in\mathrm{rng}(\mu)$.

For his specific semilattice $\mathcal{G}(\Omega)$, he proves that, whenever $|\Omega|$ is large enough,
no weakly distributive $(\vee,0)$-homomorphism
$\mu:\operatorname{Con}_c(A)\to \mathcal{G}(\Omega)$ can have $1$ in its range
(Theorem~\ref{thm:wehrung-6-1} below) \cite{Wehrung2007AdvMath}.
This nonexistence result is then used to conclude that $\mathcal{G}(\Omega)$ is not representable
as $\operatorname{Con}_c(L)$ for any lattice $L$.

\section{Wehrung's Counterexample and the Negative Solution}
\label{sec:negative}

We now sketch the mechanism behind Wehrung’s negative solution.
The proof is formulated on the level of compact congruences $\operatorname{Con}_c(L)$ and a target
distributive $(\vee,0,1)$-semilattice $\mathcal{G}(\Omega)$.

\begin{theorem}[Wehrung, Theorem 6.1]
\label{thm:wehrung-6-1}
Let $\Omega$ be a set of cardinality at least $\aleph_{\omega+1}$, and let $A$ be an algebra.
If $A$ has a congruence-compatible structure of a $(\vee,1)$-semilattice, then there is no
weakly distributive $(\vee,0)$-homomorphism
$\mu:\operatorname{Con}_c(A)\to \mathcal{G}(\Omega)$
with $1$ in its range.
\end{theorem}

\subsection*{The semilattice $\mathcal{G}(\Omega)$ construction (sketch)}
Fix a set $\Omega$.
Start with the $(\vee,0,1)$-semilattice $A(\Omega)$ generated by $1$ and elements
\[
a^0_\xi,\ a^1_\xi \qquad (\xi\in\Omega),
\]
subject to the relations $a^0_\xi\vee a^1_\xi=1$ for all $\xi\in\Omega$.
Wehrung then applies a distributive extension construction $D$ and sets
\[
\mathcal{G}(\Omega)=D(A(\Omega)).
\]
A key feature is the \emph{finite-support} property: every $x\in \mathcal{G}(\Omega)$ has a least finite support
$\operatorname{supp}(x)\subseteq\Omega$ such that $x\in \mathcal{G}(\operatorname{supp}(x))$.

\subsection*{Idea of the proof of Theorem~\ref{thm:wehrung-6-1}}
Assume toward a contradiction that $\mu:\operatorname{Con}_c(A)\to \mathcal{G}(\Omega)$ is weakly distributive and $1\in\mathrm{rng}(\mu)$.
Then:
\begin{itemize}
\item Using that $1\in\mathrm{rng}(\mu)$, pick $\alpha\in\operatorname{Con}_c(A)$ with $\mu(\alpha)=1$ and write
$\alpha$ as a finite join of principal congruences, yielding a finite decomposition
$1=\bigvee_{r<m}\mu\Theta_A(t_r,1)$.
\item For each $\xi\in\Omega$, apply weak distributivity to the decomposition $1=a^0_\xi\vee a^1_\xi$ in $\mathcal{G}(\Omega)$,
producing finite \emph{alternating} congruence chains in $A$ whose successive links map alternately below $a^0_\xi$ and $a^1_\xi$.
\item Using finite supports in $\mathcal{G}(\Omega)$ and a Kuratowski free-set argument (available at
$|\Omega|\ge\aleph_{\omega+1}$), Wehrung applies his technical lemmas (Erosion + Evaporation)
to force a collapse of these chains, contradicting the initial finite decomposition.
\end{itemize}

\hfill$\square$

For details of the construction and proof of Wehrung's theorem, see \cite[Theorem~6.1, p.~619]{Wehrung2007AdvMath}.

\subsection*{From Theorem~\ref{thm:wehrung-6-1} to the negative solution of CLP}
Let $\Omega$ be as in Theorem~\ref{thm:wehrung-6-1} and put $S:=\mathcal{G}(\Omega)$.
If $S$ were representable, i.e.\ if $S\cong\operatorname{Con}_c(K)$ for some lattice $K$,
then the isomorphism $\operatorname{Con}_c(K)\to S$ would be weakly distributive and would map some element to $1$,
contradicting Theorem~\ref{thm:wehrung-6-1}.
Hence $\mathcal{G}(\Omega)$ is not representable.
Passing from semilattices to algebraic lattices via the ideal lattice $\operatorname{Id}(\mathcal{G}(\Omega))$,
one obtains an algebraic distributive lattice that is not isomorphic to $\operatorname{Con}(K)$ for any lattice $K$,
so the Congruence Lattice Problem has a negative answer.
\\ \\
For details and other results of Wehrung's counterexample, see \cite{Wehrung2007AdvMath}.

\subsection*{R{\r u}{\v z}i{\v c}ka's optimal bound}
R{\r u}{\v z}i{\v c}ka later improved the cardinal bound to the optimal value $\aleph_2$:
there exists a distributive $(\vee,0,1)$-semilattice $G$ of size $\aleph_2$ such that no weakly distributive
$(\vee,0)$-homomorphism $\operatorname{Con}_c(A)\to G$ has $1$ in its range.
In particular, a counterexample to representability already exists at $\aleph_2$.
This bound is optimal because every distributive $(\vee,0)$-semilattice
of cardinality at most $\aleph_1$ is representable (Huhn),
so $\aleph_2$ is the smallest cardinal at which non-representability can occur. 
\\ \\ More about R{\r u}{\v z}i{\v c}ka's results in \cite{Ruzicka2008FreeTrees}


\section{Conclusion}

The Congruence Lattice Problem has a negative solution:
not every distributive algebraic lattice is isomorphic to the congruence
lattice of a lattice.
The obstruction appears only at the level of higher cardinalities,
while representability holds in all finite cases and for distributive
$(\vee,0)$-semilattices of cardinality at most $\aleph_1$.
Wehrung's construction and R{\r u}{\v z}i{\v c}ka's optimal bound
precisely determine the cardinal threshold, at which the nonrepresentability occurs.

\begin{thebibliography}{99}

\bibitem{BirkhoffFrink1948}
G.~Birkhoff and O.~Frink, Jr.,
\textit{Representations of lattices by sets},
Trans.\ Amer.\ Math.\ Soc.\ \textbf{64} (1948), 299--316.

\bibitem{FunayamaNakayama1942}
N.~Funayama and T.~Nakayama,
\textit{On the distributivity of a lattice of lattice-congruences},
Proceedings of the Imperial Academy (Tokyo) \textbf{18} (1942), 553--554.

\bibitem{GratzerGLT1998}
G.~Gr\"atzer,
\textit{General Lattice Theory},
2nd ed., new appendices by the author with B.~A.~Davey, R.~Freese, B.~Ganter,
M.~Greferath, P.~Jipsen, H.~A.~Priestley, H.~Rose, E.~T.~Schmidt, S.~E.~Schmidt,
F.~Wehrung, and R.~Wille,
Birkh\"auser Verlag, Basel, 1998, xx+663~pp.

\bibitem{Gratzer2007TwoProblems}
G.~Gr\"atzer,
\textit{Two Problems That Shaped a Century of Lattice Theory},
Notices of the American Mathematical Society \textbf{54} (2007), no.~6, 696--707.

\bibitem{GratzerCFL}
G.~Gr\"atzer,
\textit{The Congruences of a Finite Lattice: A ``Proof-by-Picture'' Approach},
3rd ed., Birkh\"auser / Springer Nature Switzerland, 2023.
ISBN 978-3-031-29062-6.
doi:\,10.1007/978-3-031-29063-3.

\bibitem{Ruzicka2008FreeTrees}
P.~R{\r u}{\v z}i{\v c}ka,
\textit{Free trees and the optimal bound in Wehrung's theorem},
Fundamenta Mathematicae \textbf{198} (2008), no.~3, 217--228.
doi:\,10.4064/fm198-3-2.

\bibitem{Wehrung2007AdvMath}
F.~Wehrung,
\textit{A solution to Dilworth's congruence lattice problem},
Advances in Mathematics \textbf{216} (2007), no.~2, 610--625.
doi:\,10.1016/j.aim.2007.05.016.

\end{thebibliography}

\end{document}



